% !TEX program = xelatex
\documentclass[11pt,a4paper,fontset=fandol]{ctexart}

% 论文常用版心：边距适中、行距略紧，避免 12pt+1.4 行距造成页面稀疏
\usepackage[margin=2.4cm,headheight=14pt]{geometry}
\usepackage{setspace}
\setstretch{1.22}

\usepackage{booktabs}
\usepackage{array}
\usepackage{tabularx}
\usepackage{amsmath,amssymb}
\usepackage{enumitem}
\usepackage{etoolbox}
\usepackage{seqsplit}
\usepackage{titlesec}
\usepackage{hyperref}
\hypersetup{hidelinks}
\usepackage{algorithm}
\usepackage{algpseudocode}
\usepackage{caption}
\usepackage{placeins}
\usepackage{flafter} % 浮动体不得出现在源码定义点之前，避免“表跑到上节正文前”
\usepackage{float}   % 仅对过高算法提供 [H] 就地放置
\usepackage{needspace} % 大块前若剩余空间不足则提前换页，避免页中部大空白

% 代码片段允许在任意字符处换行（避免长路径/标识符溢出边距）
\newcommand{\ctt}[1]{{\ttfamily\seqsplit{#1}}}
\emergencystretch=2em

% 章节标题（略收紧，减少标题前后空白）
\titleformat{\section}{\large\bfseries}{\thesection}{0.7em}{}
\titleformat{\subsection}{\normalsize\bfseries}{\thesubsection}{0.55em}{}
\titleformat{\subsubsection}{\normalsize\bfseries}{\thesubsubsection}{0.5em}{}
\titlespacing*{\section}{0pt}{1.6ex plus 0.3ex}{0.8ex}
\titlespacing*{\subsection}{0pt}{1.2ex plus 0.2ex}{0.5ex}
\titlespacing*{\subsubsection}{0pt}{1.0ex plus 0.2ex}{0.4ex}

% 图表标题：表在上、居中、小字
\captionsetup{font=small,labelsep=quad,skip=3pt}
\captionsetup[table]{position=top,justification=centering}
\captionsetup[algorithm]{font=small,labelsep=quad}

% 表格与算法统一小字号；table/algorithm 环境本身为不可拆盒子，整表不分页
\AtBeginEnvironment{table}{\footnotesize}
\AtBeginEnvironment{algorithm}{\footnotesize}

% 论文式浮动：优先顶栏；flafter 保证不倒退；floatpage 阈值略降以免孤页只放半块
\setcounter{topnumber}{3}
\setcounter{bottomnumber}{1}
\setcounter{totalnumber}{4}
\renewcommand{\topfraction}{0.9}
\renewcommand{\bottomfraction}{0.5}
\renewcommand{\textfraction}{0.1}
\renewcommand{\floatpagefraction}{0.55}
\setlength{\textfloatsep}{8pt plus 2pt minus 2pt}
\setlength{\floatsep}{6pt plus 2pt minus 2pt}
\setlength{\intextsep}{8pt plus 2pt minus 2pt}

% 公式上下间距收紧
\setlength{\abovedisplayskip}{6pt plus 2pt minus 3pt}
\setlength{\belowdisplayskip}{6pt plus 2pt minus 3pt}
\setlength{\abovedisplayshortskip}{3pt plus 1pt}
\setlength{\belowdisplayshortskip}{3pt plus 1pt}

% 段落与列表
\setlength{\parskip}{0pt plus 1pt}
\setlist[itemize]{leftmargin=1.8em,itemsep=1pt,topsep=2pt,parsep=0pt,partopsep=0pt}
\setlist[enumerate]{leftmargin=1.8em,itemsep=1pt,topsep=2pt,parsep=0pt,partopsep=0pt}

\title{\bfseries RAG 系统的评估与迭代方法\\[4pt]
\large ——以欧标（EN\,199x）中文问答系统为研究对象}
\author{Eurocode QA 评估工作组}
\date{2026 年 7 月}

\begin{document}
\maketitle

\begin{abstract}
检索增强生成（Retrieval-Augmented Generation，RAG）系统上线之后，如何在缺少大规模人工标注的条件下判断``一次改动是否带来了真实改善''，是工程迭代中的核心问题。本文以一个已投入运行的欧标（EN\,199x）中文问答系统为对象，给出一套小样本条件下可重复的评估与迭代方法。评估侧以四个核心指标刻画系统质量：忠实度（Faith）与引用精确率（CitP）度量答案生成质量并承担版本门禁，未解析引用率与上下文证据召回率作为确定性诊断信号，分别定位交叉引用补充与检索覆盖两类问题。迭代侧构成一个闭环：先由基线指标定位瓶颈环节，再依据诊断——行动矩阵选择单一变量的改动（更新参数或切换算法），通过配对 bootstrap 得到同题差值的置信区间，最终作出接受、拒绝或证据不足的三态判定。为控制评判成本，方法引入基线缓存、候选合理性前置门与检索路径差异门三道低成本闸门，使无效实验在进入昂贵的模型评判之前即被截停。本文仅描述已实现并可运行的流程；未纳入的指标及其取舍理由统一列于附录。
\end{abstract}

\clearpage
\tableofcontents

\clearpage
% ============================================================
\section{引言}
% ============================================================

\subsection{背景与问题}

面向专业规范文本的问答系统正越来越多地采用 RAG 架构\cite{rag}：先从语料库中检索与问题相关的段落，再由大语言模型依据检索结果生成带引用的回答。本文的研究对象是一个欧洲结构设计标准（EN\,199x 系列）的中文问答系统。与开放域问答不同，规范问答对答案的\emph{可追溯性}有硬性要求——每一条结论都应能落回规范条文；同时规范文本内部存在大量交叉引用（如``见 EN 1990:2002, 6.4.3.2(2)''），检索阶段能否补全这些引用直接影响答案质量。

这样的系统进入迭代期后，评估面临三重约束。其一，\textbf{人工标注稀缺}：可用的人工资产仅为客户试用期给出的 30 道题及若干自由文本评语，没有与当前系统答案一一对应的正误标签，也没有领域专家可持续参与标注。其二，\textbf{评判成本高}：以大模型充当评判者（LLM-as-judge\cite{zheng2023judge}）虽可替代人工，但每题评判都有可观的时间与费用开销，未经成本控制的评估流程会慢到无法支撑日常迭代。其三，\textbf{小样本统计}：几十道题的规模不足以宣称绝对准确率，任何版本结论都必须携带不确定性说明。

\subsection{方法概述与设计原则}

针对上述约束，本文方法遵循四条设计原则：\textbf{指标少而有用}——正文只保留能直接驱动迭代决策的四个指标，其余移入附录；\textbf{评判只用一把尺}——每题由一个与生成端不同族的评判模型一次性完成事实拆分与判定，配合内容寻址缓存，保证基线与候选版本在完全相同的标准下比较；\textbf{流程只有一条}——日常快速迭代与里程碑全量评估复用同一条管线，仅运行参数不同；\textbf{结论以可运行实现为准}——凡代码尚未实现的指标、门禁与调度机制一律不进入正文。

在这些原则之下，一轮完整的迭代包含五步：建立（或复用）基线评估；由基线指标的短板定位问题所在的流水线环节；依据诊断——行动矩阵（第\ref{sec:matrix}节）选择\emph{单一变量}的改动；在隔离环境中运行候选版本并与基线做同题配对比较；依据置信区间作出接受、拒绝或证据不足的三态判定。其中第三步的矩阵是整个闭环的枢纽：它规定了``哪个指标不达标，应当去调哪个环节的哪些参数、换哪些算法''，因此本文以最大篇幅刻画该矩阵及其使用规则。

\subsection{讨论范围}

本文不讨论需要重新训练生成模型的方案（如 Self-RAG\cite{selfrag}）、依赖开放网络搜索的检索纠错（如 CRAG\cite{crag}）、依赖外部独立知识源的事实核验（如 FActScore 的原始设定\cite{factscore}），以及需要数百条人工标注来校准评判器的方法（如 ARES 的预测驱动推断\cite{ares}）。这些方向或超出当前工程资源，或与``无人工标注''的前提冲突；相关取舍连同暂缓的指标一并记录于附录~\ref{app:deferred}。

本文其余部分组织如下：第\ref{sec:related}节回顾相关工作；第\ref{sec:system}节给出被评估系统的流水线分层，为后文的瓶颈定位建立坐标系；第\ref{sec:dataset}节描述评估数据集；第\ref{sec:metrics}节定义评估指标；第\ref{sec:iteration}节展开迭代方法论；第\ref{sec:limitations}节讨论局限；第\ref{sec:conclusion}节总结全文。

% ============================================================
\section{相关工作}
\label{sec:related}
% ============================================================

\subsection{RAG 评估框架}

近年出现了多个通用 RAG 评估框架。RAGAS\cite{ragas} 提出以忠实度（faithfulness）为核心、区分检索质量与生成质量的分层评估思想；RAGChecker\cite{ragchecker} 将答案拆解为细粒度 claim 并检查其与引用证据的对应关系；TruLens\cite{trulens} 强调``答案——上下文''支持关系的可观测性；ARES\cite{ares} 关注小样本条件下的区间估计；DeepEval\cite{deepeval} 提供测试驱动的评测组织方式。表~\ref{tab:survey}汇总了各框架对本文的启发与实际落地。本文不追求框架指标的全集，而只保留能由现有数据稳定计算、且能驱动单变量决策的部分；切块粒度、嵌入模型选择、HyDE\cite{hyde}、ColBERT\cite{colbert} 等检索侧专项改进作为候选动作纳入矩阵，但其评估仍复用同一管线。

\begin{table}[!htbp]
\centering\footnotesize
\caption{主流 RAG 评估框架的启发与本文落地}\label{tab:survey}
\begin{tabular}{@{}lll@{}}
\toprule
框架 & 保留的启发 & 本文落地 \\
\midrule
RAGAS\cite{ragas} & 忠实度与检索诊断分层 & Faith、CRec \\
RAGChecker\cite{ragchecker} & claim 与引用级判断 & Faith、CitP \\
TruLens\cite{trulens} & 答案--上下文支持关系 & Faith 的解释框架 \\
ARES\cite{ares} & 小样本区间估计 & 配对 bootstrap \\
DeepEval\cite{deepeval} & 测试驱动评测组织 & 逐题持久化与报告 \\
\bottomrule
\end{tabular}
\end{table}

\subsection{以大模型为评判者}

以大模型为评判者的可靠性与偏置已有系统研究：评判者存在位置偏好、冗长偏好等系统性偏差\cite{zheng2023judge,ye2024justice}，且当评判模型与被评模型同族时存在自我偏好\cite{wataoka2024selfpref}。结构化输出的评分方式（如 G-Eval\cite{geval}）可降低解析歧义。本文据此作出两项设计决策：评判模型与生成模型取不同模型族，以缓解自我偏好；评判结论仅用于\emph{同题配对的相对比较}而非绝对打分，使多数系统性偏差在做差时得到部分抵消（详见第\ref{sec:judge}节与第\ref{sec:limitations}节）。

\subsection{RAG 系统的参数调优}

已有工作表明，RAG 系统各环节的超参数（召回数量、重排截断、上下文预算等）对端到端质量影响显著且相互耦合\cite{wang2024best,ammar2025hyper}。这支持本文``单变量迭代''的取向：一次只改一个变量，才能把观测到的指标变化归因到具体改动上。

% ============================================================
\section{系统概述与流水线分层}
\label{sec:system}
% ============================================================

被评估系统的在线问答路径依次为：查询理解（中文问题的分解与改写）、混合检索（稠密向量检索\cite{dpr}与 BM25 稀疏检索\cite{bm25}并行召回）、融合与重排（倒数排名融合后经交叉编码器重排\cite{monobert,nogueira2019}）、上下文补充（补入被切块截断的父块，并沿规范交叉引用图扩展相关条文）、答案生成（生成带 \texttt{[Ref-$N$]} 引用标记的中文回答）。

为使``指标异常''能落到``具体环节''，本文将该路径划分为五个层级，记作 L1--L5（表~\ref{tab:stages}）。这一分层是后文瓶颈定位与诊断——行动矩阵的坐标系：每个核心指标的异常都映射到一个或两个层级，每个层级下挂一组可改动的参数与可替换的算法。

\begin{table}[!htbp]
\centering\small
\caption{在线问答流水线的五级分层}\label{tab:stages}
\begin{tabularx}{\textwidth}{@{}clXX@{}}
\toprule
层级 & 环节 & 职责 & 可观测产物 \\
\midrule
L1 & 查询理解 & 问题分解、改写与扩展 & 扩展后的查询集合 \\
L2 & 混合召回 & 向量与 BM25 并行召回候选块 & 两路召回列表 \\
L3 & 融合与重排 & 倒数排名融合、交叉编码器重排、截断 & 重排后的有序块列表 \\
L4 & 上下文补充 & 父块补全、交叉引用闭包 & 最终有序上下文、引用解析记录 \\
L5 & 答案生成 & 依据上下文生成带引用的回答 & 答案文本、\texttt{[Ref-$N$]} 标记、token 用量 \\
\bottomrule
\end{tabularx}
\end{table}

评测基础设施完整保存每题的最终答案、进入生成阶段的\emph{有序}上下文、交叉引用解析结果（已解析与未解析的集合）、token 用量与时延。这一可观测资产决定了本文方法的适用边界：它适合做同题配对的版本比较与环节级诊断，而不适合以几十道题宣称总体人群意义上的绝对准确率。

% ============================================================
\section{评估数据集}
\label{sec:dataset}
% ============================================================

\subsection{人工资产与冻结划分}

原始人工资产为客户试用期 Excel 中的 31 行记录；其中末行``其他通用问题''为界面反馈备注而非问题，予以剔除，实际入集 30 道题。每题附有正确性、完整性、多余/编造与总体判定四列自由文本评语。这些历史评语被归一化为分层字段，仅用于保证数据划分中题型与质量分布大致均衡；由于评语针对的是历史版本的答案，与当前系统答案不存在一一对应关系，本文不将其用作正误标签。

30 道题被冻结划分为开发集（15 题）、测试集（9 题）与保留集（6 题）三个互不重叠的子集。日常迭代默认在开发集上进行；测试集与保留集留待里程碑评估与最终确认，避免开发过程中的隐性过拟合。划分一经冻结不再变动，作为可复现的数据资产存在。

\subsection{参考证据（gold）的构建与使用边界}

为支持检索侧的确定性诊断，本文为每题构建参考证据：由一个模型（Codex 族）通读冻结语料，产出参考答案、原子 claims 与证据引文候选；再由另一族模型（Claude 族）对 claims 与证据的支持关系做独立复审；最后以确定性文本匹配校验每条证据引文确能在语料中逐字定位。该流程沿用原子事实核验的思想\cite{factscore,ragchecker}，以双模型交叉复审代替单模型自证。

30 道题中，15 题的参考答案被标记为\emph{语料可完整支持}（supported），9 题为\emph{部分支持}（mixed），6 题为\emph{语料缺口}（corpus\_gap，即现有语料不足以回答）。只有 supported 且证据引文非空的 15 题进入证据召回指标的分母；其余题保留在数据集中，参与不依赖参考证据的指标。需要强调的是：该参考证据未经领域专家复核，故仅用作\emph{诊断}信号，不参与版本门禁，也不能作为业务真值或上线签字依据。

\subsection{数据集的能与不能}

综上，该数据集能够回答三类问题：候选版本相对基线是否改善；问题主要出在生成环节还是检索环节；候选改动是否真正改变了检索路径。它不能给出绝对准确率结论，也不能替代结构专业人员对规范答案的审核。为此，所有评估报告均保留逐题结果、有效配对题号与不确定性区间，供人工追溯。

% ============================================================
\section{评估指标}
\label{sec:metrics}
% ============================================================

\subsection{符号与指标体系}
\label{sec:notation}

对单条样本，记 $q$ 为问题，$a$ 为系统答案，$C=(c_1,\ldots,c_m)$ 为进入生成阶段的有序上下文，$\mathcal{S}(a)$ 为从答案中拆出的原子 claims 集合，$E^+$ 为经交叉复审的证据引文集合。交叉引用解析结果分为未解析集合 $U$ 与已解析集合 $R$。版本比较总是先逐题计算指标，再在基线与候选\emph{共同}具有非空值的题上配对。

四个核心指标的分工如下：\textbf{Faith}（忠实度）是主改善指标，承担版本门禁；\textbf{CitP}（引用精确率）是非回归门，防止候选以牺牲引用纪律换取其他指标；\textbf{unresolved\_ref\_rate}（未解析引用率，下文简记 URR）与\textbf{CRec}（上下文证据召回率）是两个不依赖模型评判的确定性诊断信号，分别指向 L4 与 L1/L2 环节，不参与门禁。前两者由评判模型给出，后两者由确定性规则计算——这一分工使得``贵的指标管决策、便宜的指标管定位''。

\subsection{评判者设计}
\label{sec:judge}

Faith 与 CitP 由单一模型族的评判者给出。每题仅调用一次：评判者在一次结构化输出中完成答案的原子 claims 拆分、每条 claim 是否被上下文支持的判定，以及答案中每个 \texttt{[Ref-$N$]} 引用是否在内容上支持相邻表述的判定\cite{geval}。评判模型族与答案生成端不同族，以缓解自我偏好\cite{wataoka2024selfpref}；实际解析到的模型标识如实记入报告，不以接口名称代替模型身份。

评判失败（超时或输出不合规）的题，其 Faith 与 CitP 记为空值并计入评判失败率；当一轮评估的失败率超过 $30\%$ 时，该轮聚合结果降级为趋势参考，不用于版本判定。评判结果按内容寻址缓存：缓存键绑定模型族、实际模型标识、prompt 与输出模式版本、问题、答案及有序上下文的内容哈希。因此同一答案与上下文永不重复评判——这既控制成本，也保证基线与候选在字面意义上被``同一把尺''度量。超时、重试等实现参数见附录~\ref{app:impl}。

\subsubsection{忠实度（Faith）}

Faith 衡量答案中的原子 claims 被检索上下文支持的比例\cite{ragas,trulens}：
\begin{equation}
  \mathrm{Faith}(a,C)=\frac{1}{|\mathcal{S}(a)|}
  \sum_{s\in\mathcal{S}(a)}\mathrm{supp}(s,C),
  \qquad |\mathcal{S}(a)|>0,
  \label{eq:faith}
\end{equation}
其中 $\mathrm{supp}(s,C)\in\{0,1\}$ 表示 claim $s$ 是否被上下文 $C$ 支持。claims 为空或评判失败时该题记空值。Faith 低意味着答案在``编造''上下文没有的内容，直接指向生成环节（L5）的事实约束问题。版本判定中 Faith 采用 improve 模式：候选必须显著优于基线才可接受。

\subsubsection{引用精确率（CitP)}

CitP 衡量答案给出的引用在内容级是否成立\cite{ragchecker}：
\begin{equation}
  \mathrm{CitP}(a,C)=\frac{1}{|\mathrm{Cite}(a)|}
  \sum_{r\in\mathrm{Cite}(a)}\mathrm{valid}(r;a,C),
  \label{eq:citp}
\end{equation}
其中 $\mathrm{Cite}(a)$ 为答案中的 \texttt{[Ref-$N$]} 引用集合，$\mathrm{valid}$ 判断被引上下文是否支持引用处的相邻表述。答案非空但未给出任何引用时，CitP 记为 $0$（惩罚``有断言无出处''）；答案为空时记空值。版本判定中 CitP 采用 non-regress 模式：不要求候选提升，但显著回归即拒绝。

\subsection{确定性诊断指标}

\subsubsection{上下文证据召回率（CRec）}

CRec 回答``回答该题所需的关键证据是否被检索带回''。对每条参考证据引文做 Unicode 规范化与空白归一化后，检查其是否\emph{逐字}出现在任一上下文块中：
\begin{equation}
  \mathrm{CRec}(C,E^+)=\frac{\sum_{e\in E^+}\mathbf{1}\!\left[\mathrm{quote}(e)\subseteq C\right]}{|E^+|},
  \qquad |E^+|>0.
  \label{eq:crec}
\end{equation}
与 RAGAS 等框架用模型判断上下文相关性不同，CRec 采用确定性字符串匹配，零成本、零噪声。它只在 15 道 supported 题上计算；由于参考证据未经人工复核，CRec 仅作 L1/L2 检索覆盖的诊断信号，不参与门禁。

\subsubsection{未解析引用率（URR）}

规范文本的交叉引用能否被闭包补全是本系统特有的质量维度。设 $U$、$R$ 分别为一题中未解析与已解析的交叉引用集合，
\begin{equation}
  \mathrm{URR}=\frac{|U|}{|U|+|R|},
  \label{eq:urr}
\end{equation}
当本题不存在引用解析需求（$|U|+|R|=0$）时定义 $\mathrm{URR}=0$。URR 偏高说明上下文补充环节（L4）未能追平规范内部的引用链，是触发 L4 候选动作的前置条件（第\ref{sec:gates}节）。

\subsection{运行指标}

除上述四个质量指标外，每轮报告随附单题答案生成时延、token 用量、评判失败率与各阶段耗时。它们不参与版本判定，用于解释成本与稳定性，并为``评估流程本身是否够快''提供依据。

% ============================================================
\section{迭代方法论}
\label{sec:iteration}
% ============================================================

\subsection{总体框架}

迭代方法论回答的问题是：\emph{评估完成之后，下一步改什么、怎么确认改对了}。本文将其组织为一个五步闭环：

\begin{enumerate}
  \item \textbf{基线}：建立或复用当前版本在冻结题集上的完整评估（答案、上下文、四指标）；
  \item \textbf{定位}：按固定优先级检查基线指标，将短板映射到流水线层级（第\ref{sec:locate}节）；
  \item \textbf{选择动作}：在诊断——行动矩阵中，从该层级的候选动作里选取\emph{恰好一个}变量的改动（第\ref{sec:matrix}节）；
  \item \textbf{配对比较}：在隔离环境运行候选版本，与基线做同题配对的 bootstrap 区间估计（第\ref{sec:decide}节）；
  \item \textbf{判定与归档}：接受则候选并入新基线，拒绝则回退，证据不足则按第\ref{sec:after}节的规则处置；无论结果如何，逐题产物与判定理由均归档。
\end{enumerate}

单变量原则是整个闭环的前提：RAG 各环节超参数相互耦合\cite{wang2024best,ammar2025hyper}，一次改动多个变量将使指标变化无法归因，闭环退化为盲目搜索。

\subsection{瓶颈定位}
\label{sec:locate}

瓶颈定位以基线聚合指标为输入，按\textbf{固定优先级}依次检查，命中即停：

\begin{enumerate}
  \item $\mathrm{Faith}<0.85$：瓶颈判为 L5（生成事实约束）；
  \item 否则若 $\mathrm{CitP}<0.80$：瓶颈判为 L5（引用纪律）；
  \item 否则若 $\mathrm{URR}\ge 0.15$：瓶颈判为 L4（交叉引用补充）；
  \item 否则若 $\mathrm{CRec}<0.70$：瓶颈判为 L1/L2（检索覆盖）；
  \item 均未命中：无明确瓶颈，本轮不产生自动候选。
\end{enumerate}

将主指标置于优先级首位的理由是：Faith 是版本门禁的直接对象，若生成环节本身失真，检索侧的改善无法透过失真的生成层反映到主指标上，先行修复上游反而浪费评判预算。诊断信号（URR、CRec）排在其后，用于主指标健康时的精细定位。固定优先级还保证了定位结果可复现：同一份基线永远给出同一个瓶颈判定。

\subsection{诊断——行动矩阵}
\label{sec:matrix}

定位到层级后，``改什么''由诊断——行动矩阵（表~\ref{tab:matrix}）给出。矩阵每行对应一个流水线层级，列出触发该行的指标信号、可调整的参数类动作、可替换的算法类动作，以及各动作当前的接入状态。矩阵是迭代闭环的枢纽：它把``指标不达标''翻译成``有限的、可执行的改动清单''，后续每一轮迭代都从中取一个动作。

\begin{table}[!htbp]
\centering\footnotesize
\caption{诊断——行动矩阵：指标信号 $\to$ 流水线层级 $\to$ 候选动作}\label{tab:matrix}
\begin{tabularx}{\textwidth}{@{}llXXl@{}}
\toprule
层级 & 触发信号 & 参数类动作 & 算法类动作 & 接入状态 \\
\midrule
L1 查询理解 & CRec $<0.70$（与 L2 共享） & 多查询扩展数量；改写开关 & HyDE 式假设文档检索\cite{hyde} & 候选池 \\
L2 混合召回 & CRec $<0.70$（与 L1 共享） & 向量召回量；BM25 召回量；融合权重 & 更换嵌入模型；调整稠密/稀疏配比 & 待解锁 \\
L3 融合重排 & （无独立信号） & 重排截断数量；上下文预算 & 更换重排模型\cite{monobert}；后期交互式检索\cite{colbert} & 待解锁 \\
L4 上下文补充 & URR $\ge 0.15$ & 父块扩展深度 & 交叉引用自动闭包：关 $\to$ 开 & \textbf{已自动化} \\
L5 答案生成 & Faith $<0.85$ 或 CitP $<0.80$ & 生成温度；上下文块数预算 & 收紧生成 prompt 的事实约束与引用纪律；更换生成模型 & 候选池 \\
\bottomrule
\end{tabularx}
\end{table}

三种接入状态的含义与使用规则如下。

\textbf{已自动化}：动作已接入主循环的自动候选选择，定位命中即可无人工介入地完成整轮比较。当前唯一经过真实数据流验证的自动动作是 L4 的交叉引用自动闭包开关——已验证该开关确实改变检索输出，且可经由隔离环境的配置注入生效。

\textbf{待解锁}：动作语义明确，但当前主路径存在初始截断常量的钳制，使 L2/L3 的召回量与重排截断参数在一定范围内\emph{不改变最终输出}。在该钳制解除并验证旋钮确实产生输出差异之前，将这些动作纳入自动选择只会浪费评估预算，故暂列为待解锁。这一状态本身即是矩阵的产出之一：它把``调了没用的参数''显式登记为工程债，而不是留给试错。

\textbf{候选池}：动作真实有效但需人工实施（如修改生成 prompt、更换模型），实施后仍走完全相同的配对比较管线获得三态判定。换言之，矩阵不要求所有动作自动化，但要求\emph{所有动作的效果确认方式一致}。

两点补充。其一，CRec 无法区分 L1 与 L2：证据未被带回可能因为查询没问对，也可能因为召回没捞到。区分依靠一次人工检视——检查扩展后的查询集合是否覆盖了证据所在条文的术语；这是当前流程中唯一需要人工判断的定位步骤。其二，L3 暂无独立信号：当 CRec 正常而 Faith 偏低、且逐题检视发现关键证据被排序挤出上下文窗口时，应怀疑重排环节；建立 L3 的独立信号（如证据在重排后的排名统计）列为后续工作。

\subsection{成本闸门}
\label{sec:gates}

评判是流程中最昂贵的操作，因此主循环在评判之前设置三道低成本闸门，使注定无效的实验尽早终止：

\textbf{基线缓存}。基线评估按代码版本、基线配置与数据集版本三元组缓存；三者均未变化时直接复用既有基线，不重复生成与评判。强制重跑仅在里程碑档位使用。

\textbf{候选合理性门}。自动候选须与瓶颈定位一致方可运行：例如 L4 动作要求基线 $\mathrm{URR}\ge 0.15$。定位指向其他层级而无相应自动动作时，本轮直接以``证据不足''结束，不产生任何生成与评判开销——把``方向不对''的结论从一小时压缩到一秒。

\textbf{检索路径差异门}。候选答案生成完成后、进入评判之前，逐题比较基线与候选的有序上下文块序列。发生变化的题目不足 3 题时，判定候选未真正改变检索路径，本轮以``证据不足''结束并跳过评判——配置开关未生效、改动被上游钳制吸收等常见的``假候选''在此被拦截。序列比较是纯内存操作，成本可忽略；顺序变化亦计入差异，因为顺序影响生成。

\subsection{主循环}

算法~\ref{alg:loop}给出主循环的完整控制流，其中候选答案逐题落盘、支持断点续跑，单题失败不终止整轮（实现细节见附录~\ref{app:impl}）。

\begin{algorithm}[H]
\caption{单管线三态比较主循环}\label{alg:loop}
\begin{algorithmic}[1]
\Require 冻结题集 $D$，有效配对下限 $\mathrm{gate\_n}=12$
\State $B\gets$\Call{LoadOrRunBaseline}{$D$} \Comment{按代码版本、配置、数据版本缓存}
\State $L\gets$\Call{LocateBottleneck}{$B$} \Comment{固定优先级：Faith $\to$ CitP $\to$ URR $\to$ CRec}
\State $a\gets$\Call{PickAction}{$L$} \Comment{诊断——行动矩阵中该层级的自动动作}
\If{$a$ 为空，或 $a$ 的前置信号未达阈值}
  \State \Return inconclusive \Comment{候选合理性门；不产生任何评估开销}
\EndIf
\State $A_C\gets$\Call{RunCandidateAnswers}{$a,D$} \Comment{隔离环境；逐题落盘、断点续跑}
\State $A_B\gets$\Call{LoadAnswers}{$B$}
\State $K\gets$\Call{OrderedContextDiff}{$A_B,A_C$}
\If{$|K|<3$}
  \State \Return inconclusive \Comment{检索路径差异门；跳过候选评判}
\EndIf
\State $J_C\gets$\Call{JudgeCandidate}{$A_C$} \Comment{基线评判结果自缓存读取}
\State \Return \Call{ABDecide}{$B,J_C,\mathrm{gate\_n}$} \Comment{三态判定，见第\ref{sec:decide}节}
\end{algorithmic}
\end{algorithm}

\subsection{配对 bootstrap 与三态判定}
\label{sec:decide}

设基线与候选在第 $i$ 个共同有效题上的指标值分别为 $x_i$ 与 $y_i$，逐题差值 $d_i=y_i-x_i$。以固定随机种子对 $\{d_i\}$ 做 $B=1000$ 次有放回配对重采样，取均值差的经验 $95\%$ 置信区间 $[d_L,d_U]$\cite{ares}。判定阈值取 $\epsilon=\delta=0.05$，有效配对下限 $\mathrm{gate\_n}=12$，区间宽度上限 $0.1$：

\begin{itemize}[leftmargin=2.5em]
  \item 有效配对数 $n<\mathrm{gate\_n}$：inconclusive（样本不足以支撑任何方向的结论）；
  \item 任一指标满足 $d_U<-\delta$：reject（存在显著回归）；
  \item Faith（improve 模式）：$d_L>+\epsilon$ 且区间宽度不超过 $0.1$ 时 accept，否则 inconclusive；
  \item CitP（non-regress 模式）：未见显著回归且区间宽度不超过 $0.1$ 时通过；宽度超限时 inconclusive；
  \item 合并规则：任一指标 reject 则整体 reject；Faith accept 且 CitP 门通过才整体 accept；其余情形均为 inconclusive。
\end{itemize}

区间宽度上限的作用是防止``碰巧显著''：区间过宽说明逐题差值高度不稳定，即便端点越过阈值也只作探索性结论。评判失败率超过 $30\%$ 的轮次整体降级为趋势参考，不进入上述判定。

\subsection{判定后的动作}
\label{sec:after}

\textbf{accept}：候选改动并入代码或基线配置，代码版本随之变化，下一轮自动以新版本重建基线；旧基线及其判定报告归档留痕。\textbf{reject}：改动回退，判定理由（哪个指标、区间端点）写入报告，该动作在矩阵中标记``已试、显著回归''，避免重复试探。\textbf{inconclusive}：按触发原因分流——因样本不足触发的，可扩大划分（如从开发集转向全量）后重试同一动作；因差异门触发的，说明动作未生效，应先排查配置注入或上游钳制；因区间过宽触发的，记为探索性结论，供后续合并证据。三种出口都不允许``悄悄重跑直到显著''：每轮判定连同随机种子与逐题差值一并归档，重复运行同一比较得到同一结论。

\subsection{运行档位}

日常与里程碑复用同一条管线。\textbf{日常档}在开发集上运行，复用基线缓存，一轮的边际成本主要是候选答案生成与增量评判。\textbf{里程碑档}在全量题集上强制重建基线并重复两次：两次基线之间的指标波动作为自噪声（self-noise）写入报告，为解读候选差值提供参照——若候选带来的差值与自噪声同量级，即便区间显著也应谨慎解读。自噪声仅作展示，不改变判定阈值。

% ============================================================
\section{局限与讨论}
\label{sec:limitations}
% ============================================================

\textbf{单评判者的偏差}。单一评判模型无法消除位置偏置、风格偏好与绝对分数的校准误差\cite{zheng2023judge,ye2024justice}。本文以三重手段控制风险：评判者与生成端不同族；内容寻址缓存保证基线与候选使用同一把尺；结论只取同题差值与置信区间。因此方向性结论（候选是否优于基线）可用于迭代决策，绝对分数（如 Faith$=0.85$ 的含义）仍应谨慎解读。

\textbf{参考证据未经人工复核}。gold 由双模型交叉复审产生，可能存在两模型共同的盲区。本文将其严格限定在诊断用途，任何版本门禁都不依赖 gold；一旦获得领域专家复核，附录~\ref{app:deferred}中依赖 gold 的指标（如答案完整性）可按原定义启用。

\textbf{小样本}。30 题的规模决定了本方法只能支撑相对比较。有效配对下限与区间宽度上限把``样本不足''显式转化为 inconclusive，而非给出虚假的确定性；代价是细小但真实的改善可能需要多轮积累证据才能被接受。

\textbf{矩阵的覆盖缺口}。L3 缺少独立诊断信号，L1 与 L2 共享信号需一次人工区分；L2/L3 的参数动作受上游钳制待解锁。矩阵的价值在于把这些缺口显式登记为工程债并给出解锁条件，而非掩盖它们。

% ============================================================
\section{结论}
\label{sec:conclusion}
% ============================================================

本文面向一个已上线的欧标中文 RAG 问答系统，给出了小样本、无持续人工标注条件下的评估与迭代方法：四个各司其职的核心指标（门禁两个、诊断两个），一个与生成端不同族、带内容寻址缓存的单评判者，以及以诊断——行动矩阵为枢纽的单变量迭代闭环。三道成本闸门使无效实验在进入昂贵评判前即被截停，配对 bootstrap 的三态判定则把小样本的不确定性显式暴露在每一次版本决策中。方法的全部构件均有可运行实现支撑；后续工作按矩阵登记的缺口推进：解除召回截断钳制以解锁 L2/L3 参数动作、建立 L3 的独立诊断信号，以及在获得专家复核后启用依赖参考证据的门禁指标。

\clearpage
% ============================================================
\begin{thebibliography}{99}

\bibitem{rag}
P.~Lewis, E.~Perez, A.~Piktus, et al.
\emph{Retrieval-Augmented Generation for Knowledge-Intensive NLP Tasks.}
NeurIPS 2020. arXiv:2005.11401.
\url{https://arxiv.org/abs/2005.11401}

\bibitem{dpr}
V.~Karpukhin, B.~Oguz, S.~Min, et al.
\emph{Dense Passage Retrieval for Open-Domain Question Answering.}
EMNLP 2020. arXiv:2004.04906.
\url{https://arxiv.org/abs/2004.04906}

\bibitem{bm25}
S.~Robertson, H.~Zaragoza.
\emph{The Probabilistic Relevance Framework: BM25 and Beyond.}
Foundations and Trends in Information Retrieval, 2009.

\bibitem{monobert}
R.~Nogueira, K.~Cho.
\emph{Passage Re-ranking with BERT.}
arXiv:1901.04085, 2019.
\url{https://arxiv.org/abs/1901.04085}

\bibitem{nogueira2019}
R.~Nogueira, W.~Yang, K.~Cho, J.~Lin.
\emph{Multi-Stage Document Ranking with BERT.}
arXiv:1910.14424, 2019.
\url{https://arxiv.org/abs/1910.14424}

\bibitem{colbert}
O.~Khattab, M.~Zaharia.
\emph{ColBERT: Efficient and Effective Passage Search via Contextualized Late Interaction over BERT.}
SIGIR 2020. arXiv:2004.12832.
\url{https://arxiv.org/abs/2004.12832}

\bibitem{hyde}
L.~Gao, X.~Ma, J.~Lin, J.~Callan.
\emph{Precise Zero-Shot Dense Retrieval without Relevance Labels.}
ACL 2023. arXiv:2212.10496.
\url{https://arxiv.org/abs/2212.10496}

\bibitem{ragas}
S.~Es, J.~James, L.~Espinosa-Anke, S.~Schockaert.
\emph{RAGAS: Automated Evaluation of Retrieval Augmented Generation.}
arXiv:2309.15217, 2023.
\url{https://arxiv.org/abs/2309.15217}

\bibitem{ares}
J.~Saad-Falcon, O.~Khattab, C.~Potts, M.~Zaharia.
\emph{ARES: An Automated Evaluation Framework for Retrieval-Augmented Generation Systems.}
NAACL 2024. arXiv:2311.09476.
\url{https://arxiv.org/abs/2311.09476}

\bibitem{ragchecker}
D.~Ru, J.~Qiu, L.~Hu, et al.
\emph{RAGChecker: A Fine-grained Framework for Diagnosing Retrieval-Augmented Generation.}
NeurIPS 2024 (Datasets and Benchmarks). arXiv:2408.08067.
\url{https://arxiv.org/abs/2408.08067}

\bibitem{deepeval}
Confident AI.
\emph{DeepEval: The Open-Source Evaluation Framework for LLMs.}
\url{https://github.com/confident-ai/deepeval}, 2023--2025.

\bibitem{geval}
Y.~Liu, D.~Iter, Y.~Xu, et al.
\emph{G-Eval: NLG Evaluation using GPT-4 with Better Human Alignment.}
arXiv:2303.16634, 2023.
\url{https://arxiv.org/abs/2303.16634}

\bibitem{trulens}
TruEra / Snowflake.
\emph{TruLens: LLM App Observability \& Evaluation.}
\url{https://www.trulens.org/}, 2023--2025.

\bibitem{selfrag}
A.~Asai, Z.~Wu, Y.~Wang, et al.
\emph{Self-RAG: Learning to Retrieve, Generate, and Critique through Self-Reflection.}
ICLR 2024. arXiv:2310.11511.
\url{https://arxiv.org/abs/2310.11511}

\bibitem{crag}
S.~Yan, J.~Gu, Y.~Zhu, Z.~Ling.
\emph{Co{}rrective Retrieval Augmented Generation.}
arXiv:2401.15884, 2024.
\url{https://arxiv.org/abs/2401.15884}

\bibitem{factscore}
S.~Min, K.~Krishna, J.~Lyu, et al.
\emph{FActScore: Fine-grained Atomic Evaluation of Factual Precision in Long Form Text Generation.}
EMNLP 2023. arXiv:2305.14251.
\url{https://arxiv.org/abs/2305.14251}

\bibitem{wang2024best}
X.~Wang, Z.~Wang, X.~Gao, et al.
\emph{Searching for Best Practices in Retrieval-Augmented Generation.}
EMNLP 2024. arXiv:2407.01219.
\url{https://arxiv.org/abs/2407.01219}

\bibitem{ammar2025hyper}
M.~Ammar, et al.
\emph{Optimizing Retrieval-Augmented Generation: Analysis of Hyperparameter Impact on Performance and Efficiency.}
arXiv:2505.08445, 2025.
\url{https://arxiv.org/abs/2505.08445}

\bibitem{zheng2023judge}
L.~Zheng, W.~Chiang, Y.~Sheng, et al.
\emph{Judging LLM-as-a-Judge with MT-Bench and Chatbot Arena.}
NeurIPS 2023 (Datasets and Benchmarks). arXiv:2306.05685.
\url{https://arxiv.org/abs/2306.05685}

\bibitem{wataoka2024selfpref}
R.~Wataoka, M.~Takahashi, K.~Ri.
\emph{Self-Preference Bias in LLM-as-a-Judge.}
arXiv:2410.21819, 2024.
\url{https://arxiv.org/abs/2410.21819}

\bibitem{ye2024justice}
Z.~Ye, X.~Wang, Y.~Huang, et al.
\emph{Justice or Prejudice? Quantifying Biases in LLM-as-a-Judge.}
arXiv:2410.02736, 2024.
\url{https://arxiv.org/abs/2410.02736}

\end{thebibliography}

\clearpage
\appendix

% ============================================================
\section{实现细节}
\label{app:impl}
% ============================================================

本附录汇总正文有意省略的实现参数与工程约定，全部与 \ctt{eval\_methodology/mvp/} 下的代码常量一致。

\subsection{隔离与候选注入}

基线与候选均以独立子进程（sidecar）运行被评估服务，候选配置仅以环境变量注入子进程，主项目的 \texttt{.env} 不被写入。每轮运行前后记录 git 脏文件集合与 \texttt{.env} 的 SHA-256：若出现 \ctt{eval\_methodology/} 之外的新脏路径或环境文件变化，判定隔离失败并在报告中标记。

\subsection{答案资产与断点续跑}

答案按运行标识、版本（baseline/candidate）与题号分层存放于 \ctt{artifacts/runs/<run\_id>/answers/<variant>/<qid>.json}。每题完成即以临时文件原子替换写入；断点续跑时跳过状态为成功的题，单题失败记录错误类型与信息，整轮继续。缓存与续跑均以这些文件为唯一事实来源。

\subsection{评判者实现参数}

评判通过命令行客户端请求符合 JSON Schema 的结构化输出，默认使用 Claude 族，可经环境变量 \texttt{MVP\_JUDGE\_FAMILY} 切换为 Codex 族。单次调用超时 180 秒，失败后重试 1 次。prompt 仅包含前 30 个有序上下文块，每块截断至 800 字符；客户端禁用工具调用。缓存键为内容寻址哈希，绑定模型族、实际解析到的模型标识、prompt/schema 版本、问题、答案与有序上下文内容。

\subsection{基线缓存与判定常量}

基线缓存键由 \texttt{git rev-parse HEAD}、基线配置哈希与数据集版本号三部分组成；\texttt{--fresh} 强制重建。判定常量：$B=1000$（固定种子）、$\epsilon=\delta=0.05$、$\mathrm{gate\_n}=12$、区间宽度上限 $0.1$、评判失败率降级阈值 $30\%$、检索路径差异门 3 题、URR 前置门 $0.15$、证据引文最短长度 20 字符。

\subsection{报告与运行命令}

JSON 与 Markdown 报告同时包含：逐题四指标与状态、聚合值与版本差值、三态判定及理由、有效配对题号、\ctt{answers/judge/decide/report} 四阶段耗时、实际评判模型标识、平均时延、token 合计与评判失败率。常用命令：

\begin{verbatim}
# 日常：基线（开发集，缓存复用）
uv run python -m eval_methodology.mvp.loop.run_mvp baseline --split dev
# 日常：单变量比较
uv run python -m eval_methodology.mvp.loop.run_mvp compare --split dev
# 里程碑：全量、强制重建、基线重复两次
uv run python -m eval_methodology.mvp.loop.run_mvp baseline \
    --split full --repeats 2 --fresh
# 文档编译
cd eval_methodology && latexmk -xelatex -interaction=nonstopmode main.tex
\end{verbatim}

% ============================================================
\section{指标可算性汇总}
% ============================================================

\begin{table}[!htbp]
\centering\footnotesize
\caption{指标、数据依赖与角色}\label{tab:feasibility}
\begin{tabularx}{\textwidth}{@{}lXXl@{}}
\toprule
指标 & 数据依赖 & 当前口径 & 角色 \\
\midrule
Faith & 答案、有序上下文 & 单评判者 claims 支持率 & 主指标/门禁 \\
CitP & 答案引用、有序上下文 & 单评判者引用有效率 & 非回归门 \\
URR & 引用解析记录 & 确定性比例 & L4 诊断 \\
CRec & supported gold 引文、上下文 & 归一化引文精确覆盖 & L1/L2 诊断 \\
单题时延 & 答案检查点 & 毫秒 & 随报 \\
token 用量 & 生成端用量记录 & 总量汇总 & 随报 \\
评判失败率 & 评判调用状态 & 失败题/总题数 & 随报/趋势降级 \\
阶段耗时 & 四阶段计时 & 秒 & 随报 \\
\bottomrule
\end{tabularx}
\end{table}

% ============================================================
\section{延后指标清单}
\label{app:deferred}
% ============================================================

\begin{table}[!htbp]
\centering\footnotesize
\caption{未进入正文门禁的历史指标及延后理由}\label{tab:deferred}
\begin{tabularx}{\textwidth}{@{}lXX@{}}
\toprule
指标 & 一句话定义 & 延后理由 \\
\midrule
Hall & $1-\mathrm{Faith}$ 的镜像 & 与 Faith 重复，不单独计数 \\
Match & 评判总体等级与人工等级的一致率 & 历史人工标签不对应当前答案 \\
Comp & 答案覆盖 gold claims 的比例 & 依赖 gold claims 与评判族独立性 \\
Corr & 答案相对参考答案的正确性 & 参考答案未经人工复核，不足以作门禁 \\
CPrec & 有序上下文中 gold 证据的排名精确率 & 证据密度与等价证据未标定 \\
Noise & 负证据进入上下文的比例 & 需要稳定的负证据标注 \\
ECov & 进入证据阶段的 gold 证据覆盖率 & 需要单独的证据阶段集合 \\
RDeg & 重排失败并回退的比例 & 需修改主项目埋点，超出隔离边界 \\
$\tau$ & 系统置信等级与 Faith 的 Kendall 相关 & 三档置信度并列值过多 \\
\bottomrule
\end{tabularx}
\end{table}

% ============================================================
\section{交叉审阅摘要}
\label{sec:trail}
% ============================================================

正文初稿由一个模型起草，由另一模型（Codex）做对抗审阅后逐条裁定修改。首轮共 20 条意见（高 8 / 中 9 / 低 3），全部采纳，修改集中在伪代码可执行性、小样本统计前提、杠杆分类与引用表述一致性。v2 进一步按``可运行优先''原则收敛：正文只保留已实现的四个核心指标、单模型族评判与单管线三态决策；收敛决策记录于 \ctt{MVP\_V2\_IMPL\_PLAN.md}，实现入口见 \ctt{mvp/README.md}。v3（本稿）按标准论文结构重组：新增系统分层与结论章节，迭代方法论扩展为诊断——行动矩阵，实现参数下沉至附录~\ref{app:impl}。

\end{document}
